% GENERATED FILE. Edit canonical lessons, then run npm run generate:books.
% canonical-source-hash: fnv1a64:f48445f364a7dea8
% canonical-lessons: GE-C168-unmoeglich, GE-C168-hell, GE-C168-dunkel, GE-C168-voll, GE-C168-altmodisch

\chapter{Impossible, Bright, Dark, Full, Old-fashioned}
\label{ch:ge-impossible-bright-dark-full-old-fashioned}
\hlchaptermodality{168}

\begin{chapteropening}
\textbf{By the end of this chapter:} \emph{I can say impossible, bright, dark, full and old-fashioned.}

Complete the last lesson of chapter 168: I can say impossible, bright, dark, full and old-fashioned.
\end{chapteropening}

\section[unmöglich]{unmöglich — impossible}

\label{lesson:GE-C168-unmoeglich}



\begin{quote}
Before the new one: say the German for exciting, then the German for possible.
\end{quote}

\subsection*{You'll want to know: unmöglich}
\textbf{unmöglich} — \textquotedblleft{}impossible\textquotedblright{}.

\begin{cousinweb}
\textbf{unmöglich} is \emph{un-}, \textquotedblleft{}not\textquotedblright{}, + \emph{möglich}, \textquotedblleft{}possible\textquotedblright{}, which is built on \emph{mögen}, the verb behind English \emph{may}. (EtymDB 2.1)
\end{cousinweb}

\begin{grammarlens}[title={What you've built}]
One more word for this chapter.
\end{grammarlens}

\subsection*{Your turn}
\begin{itemize}
\raggedright
  \item \emph{Say it:} \emph{unmöglich}
  \item \emph{Say it:} \emph{unmöglich}, once more
  \item \emph{Say it:} say \emph{unmöglich}
  \item \emph{Recall:} say the German for exciting, then the German for possible, then say \emph{unmöglich} again
\end{itemize}

\begin{tcolorbox}[breakable,colback=teal!4,colframe=teal!35!black,title={Before you move on}]
What does unmöglich mean? (\textquotedblleft{}Impossible\textquotedblright{}.) Say it once more.
\end{tcolorbox}
\section[hell]{hell — bright, light}

\label{lesson:GE-C168-hell}



\begin{quote}
Before the new one: say the German for possible, then the German for impossible.
\end{quote}

\subsection*{You'll want to know: hell}
\textbf{hell} — \textquotedblleft{}bright, light\textquotedblright{}.

\begin{cousinweb}
\textbf{hell} $\leftarrow$ Middle High German \emph{hel}, \textquotedblleft{}ringing, clear-sounding\textquotedblright{} — inherited; EtymDB 2.1 lists Dutch \emph{hel} beside it. A word for a clear sound came to mean a clear light.
\end{cousinweb}

\begin{grammarlens}[title={What you've built}]
One more word for this chapter.
\end{grammarlens}

\subsection*{Your turn}
\begin{itemize}
\raggedright
  \item \emph{Say it:} \emph{hell}
  \item \emph{Say it:} \emph{hell}, once more
  \item \emph{Say it:} say \emph{hell}
  \item \emph{Recall:} say the German for possible, then the German for impossible, then say \emph{hell} again
\end{itemize}

\begin{tcolorbox}[breakable,colback=teal!4,colframe=teal!35!black,title={Before you move on}]
What does hell mean? (\textquotedblleft{}Bright, light\textquotedblright{}.) Say it once more.
\end{tcolorbox}
\section[dunkel]{dunkel — dark}

\label{lesson:GE-C168-dunkel}



\begin{quote}
Before the new one: say the German for impossible, then the German for bright.
\end{quote}

\subsection*{You'll want to know: dunkel}
\textbf{dunkel} — \textquotedblleft{}dark\textquotedblright{}.

\begin{cousinweb}
\textbf{dunkel} $\leftarrow$ Middle High German \emph{tunkel} $\leftarrow$ Old High German \emph{tunkal}, \textquotedblleft{}dark\textquotedblright{} (EtymDB 2.1) — inherited.
\end{cousinweb}

\begin{grammarlens}[title={What you've built}]
One more word for this chapter.
\end{grammarlens}

\subsection*{Your turn}
\begin{itemize}
\raggedright
  \item \emph{Say it:} \emph{dunkel}
  \item \emph{Say it:} \emph{dunkel}, once more
  \item \emph{Say it:} say \emph{dunkel}
  \item \emph{Recall:} say the German for impossible, then the German for bright, then say \emph{dunkel} again
\end{itemize}

\begin{tcolorbox}[breakable,colback=teal!4,colframe=teal!35!black,title={Before you move on}]
What does dunkel mean? (\textquotedblleft{}Dark\textquotedblright{}.) Say it once more.
\end{tcolorbox}
\section[voll]{voll — full}

\label{lesson:GE-C168-voll}



\begin{quote}
Before the new one: say the German for bright, then the German for dark.
\end{quote}

\subsection*{You'll want to know: voll}
\textbf{voll} — \textquotedblleft{}full\textquotedblright{}.

\begin{cousinweb}
\textbf{voll} $\leftarrow$ Middle High German \emph{vol} $\leftarrow$ Old High German \emph{foll}, \textquotedblleft{}full\textquotedblright{} — inherited, the same word as English \emph{full}. (EtymDB 2.1)
\end{cousinweb}

\begin{grammarlens}[title={What you've built}]
One more word for this chapter.
\end{grammarlens}

\subsection*{Your turn}
\begin{itemize}
\raggedright
  \item \emph{Say it:} \emph{voll}
  \item \emph{Say it:} \emph{voll}, once more
  \item \emph{Say it:} say \emph{voll}
  \item \emph{Recall:} say the German for bright, then the German for dark, then say \emph{voll} again
\end{itemize}

\begin{tcolorbox}[breakable,colback=teal!4,colframe=teal!35!black,title={Before you move on}]
What does voll mean? (\textquotedblleft{}Full\textquotedblright{}.) Say it once more.
\end{tcolorbox}
\section[altmodisch]{altmodisch — old-fashioned}

\label{lesson:GE-C168-altmodisch}



\begin{quote}
Before the new one: say the German for dark, then the German for full.
\end{quote}

\subsection*{You'll want to know: altmodisch}
\textbf{altmodisch} — \textquotedblleft{}old-fashioned\textquotedblright{}.

\begin{cousinweb}
\textbf{altmodisch} is \emph{alt}, \textquotedblleft{}old\textquotedblright{}, + the word for \textquotedblleft{}fashion\textquotedblright{}, from French \emph{mode} $\leftarrow$ Latin \emph{modus}, \textquotedblleft{}manner\textquotedblright{}, + \emph{-isch}: in the old manner.
\end{cousinweb}

\begin{grammarlens}[title={What you've built}]
One more word for this chapter.
\end{grammarlens}

\subsection*{Your turn}
\begin{itemize}
\raggedright
  \item \emph{Say it:} \emph{altmodisch}
  \item \emph{Say it:} \emph{altmodisch}, once more
  \item \emph{Say it:} say \emph{altmodisch}
  \item \emph{Recall:} say the German for dark, then the German for full, then say \emph{altmodisch} again
\end{itemize}

\begin{tcolorbox}[breakable,colback=teal!4,colframe=teal!35!black,title={Before you move on}]
What does altmodisch mean? (\textquotedblleft{}Old-fashioned\textquotedblright{}.) Say it once more.
\end{tcolorbox}
